\chapter{Training/inference parity and parallel recurrence}
\label{evod26:chapter}
A recurrent architecture is defined by its transition and state boundaries, not by
a particular Python loop. If a different execution schedule implements the same
function, training can exploit parallel work while inference updates one step at
a time. The difficult part is proving equivalence without overlooking gradients,
resets, padding, floating-point order or request ownership.

Chapter 25 tested actual Chapter 22 DOGMA carry and actual Chapter 23 Hermon causal
prefix behavior. Here we isolate an affine recurrent core and build four independent
views of it: a literal sequential loop, an inclusive scan, a chunked loop and a
closed-form product/sum oracle. We then derive a reverse-mode gradient oracle.
These are algebra and software experiments, not trained genomic capability or
accelerator performance evidence.

\section{Specify the transition before changing the schedule}
Consider
\begin{equation}
h_t=a_t\odot h_{t-1}+b_t,\qquad h_t\in\mathbb R^D,\quad t=1,\ldots,T.
\label{evod26:rec}
\end{equation}
The elementwise coefficients $a_t,b_t$ can be learned or input-dependent.
For a whole-sequence parallel scan they must be available without first computing
the recurrent states that the scan is supposed to produce. Causal legality alone
is not sufficient: $a_t=g(h_{t-1},x_t)$ may be causal yet require sequential
evaluation. Local input-derived coefficients are one scan-compatible case.

This distinction matters for DOGMA. Its Chapter 22 complete state contains native
hidden state and memory loci, including state-dependent read/write interactions.
The entire nonlinear state machine does not become scan-compatible merely because
one subexpression is affine. A scan proof for a frozen or input-conditioned core
does not prove a scan implementation of that full model. HermonDNA remains the
Transformer branch; Evolutor coordinates both rather than conflating their state.
\Plate{01-three}{One affine transition has sequential, chunked and scan schedules.
Coefficient availability is a prerequisite; this is not a claim that every DOGMA
transition admits the same scan.}{evod26:three}

The reference uses $a,b$ shaped $[T,*\mathrm{state}]$ and $h_0$ shaped
$[*\mathrm{state}]$. Extra state axes can hold independent batch lanes.
Shapes, floating dtype, device and finite inputs must match. It does not mutate
PyTorch's global default dtype. The empty sequence returns no per-position states
and final state $h_0$, rather than inventing a transition.

\section{Derive chronological affine composition}
Let $F=(a,b)$ denote $F(h)=a\odot h+b$. Applying $F_L$ then $F_R$ gives
\[
F_R(F_L(h))=(a_R\odot a_L)\odot h+(a_R\odot b_L+b_R).
\]
Define chronological composition
\begin{equation}
F_L\otimes F_R=(a_R\odot a_L,\ a_R\odot b_L+b_R).
\label{evod26:composition}
\end{equation}
The left operand happens earlier even though function composition is written
$F_R\circ F_L$. This convention must remain consistent in the diagram, code and tests.
\Plate{02-compose}{Chronological order fixes the translated term $a_Rb_L+b_R$.
Associativity permits regrouping; it never permits reordering.}{evod26:compose}
\Listing{compose}

For three steps both parenthesizations produce
\[
(a_3a_2a_1,\ a_3a_2b_1+a_3b_2+b_3),
\]
componentwise. Thus the pairs with identity $(\mathbf1,\mathbf0)$ form an
associative monoid over exact real arithmetic. They are generally not commutative.
For $F_1=(2,1)$, $F_2=(3,4)$, chronological $F_1\otimes F_2=(6,7)$,
whereas $F_2\otimes F_1=(6,9)$. A parallel algorithm that swaps operands is wrong.

An independent test represents scalar maps by homogeneous matrices
$\begin{pmatrix}a&b\\0&1\end{pmatrix}$. Their multiplication supplies an oracle
without calling \texttt{compose}. Matrix-affine recurrences also compose in
principle, but their operation cost and coefficient structure differ from this
elementwise implementation; the scalar cost model must not be transferred blindly.

\section{Prefixes, not just one tree reduction}
The \Term{affine prefix summary} at position $t$ is
\[
P_t=F_1\otimes\cdots\otimes F_t=(A_t,B_t),\qquad
h_t=A_t\odot h_0+B_t.
\]
Training often needs every $h_t$, not just the terminal state. A balanced reduction
that returns only $P_T$ is therefore not the required algorithm.

Our transparent Hillis--Steele scan starts with one-step pairs. At offsets
$1,2,4,\ldots$, a position combines its earlier interval with its current interval
using arrays from the \emph{previous} stage. Unaffected leading positions are copied.
After stage $k$, each position summarizes up to $2^{k+1}$ preceding steps including
itself. This invariant proves complete prefixes after $\lceil\log_2T\rceil$ stages.
\Plate{03-tree}{An inclusive scan produces every prefix. At each doubling stage
all updates read the preceding stage; in-place partial updates would violate the
interval invariant.}{evod26:tree}
\Listing{hillis_steele_scan}

For $h_0=2$, $a=(0.5,2,0,-1)$ and $b=(1,0.5,3,2)$, direct updates give
$h=(2,4.5,3,-1)$. The prefix summaries expose the same result:
\begin{center}\input{\Generated/trace.tex}\end{center}
A zero coefficient discards the incoming state. Negative and expansive coefficients
are legal algebraically; imposing $(0,1)$ coefficients in a model is an additional
parameterization choice, not a requirement of affine composition.

Blelloch's chapter establishes associative prefix computation as prior art
\cite{Blelloch1993}. S5's author abstract describes a multi-input, multi-output
state-space layer compatible with parallel scans \cite{Smith2023}. Mamba's selected
mechanism sections distinguish input-dependent state-space computation and
hardware-aware scan, fusion and recomputation \cite{GuDao2024}. These sources
motivate engineering questions; our reference neither reproduces their trained
results nor implements their production kernels.

\section{Resets and padding belong in the algebra}
Let $r_t$ mark \Term{reset-before-update} and $v_t$ mark a valid position. With an explicit
reset state $h_{\rm reset}$, a valid reset step computes
$a_t\odot h_{\rm reset}+b_t$ independently of $h_{t-1}$.
Fold these semantics into effective pairs:
\[
(\tilde a_t,\tilde b_t)=
\begin{cases}
(\mathbf1,\mathbf0),&v_t=0,\\
(\mathbf0,a_t\odot h_{\rm reset}+b_t),&v_t=1,\ r_t=1,\\
(a_t,b_t),&v_t=1,\ r_t=0.
\end{cases}
\]
Now the same scan machinery handles segmented records and identity padding.
Our schema rejects reset requests on padding rather than allowing conflicting
instructions. Masks are Boolean time vectors shared across state lanes; per-lane
ragged resets require a separately extended mask contract, not accidental broadcasting.
\Listing{effective_maps}

For $h_0=2$, $h_{\rm reset}=10$, $a_t=0.5$, $b_t=1$, three valid steps with
a reset at step two give $(2,6,4)$. If step three is padding, its state remains $6$.
Setting padded $a=b=0$ would erase the state, not preserve it.
A reset cuts dependence on the earlier record, but may preserve dependence on a
shared learned reset parameter. When $h_{\rm reset}=h_0$, gradients to that common
parameter still accumulate across reset segments.

\section{Carry the state and preserve its gradient}
A chunk boundary is not a biological sequence boundary. For a chunk starting at
$s$, its incoming state must be the previous chunk's terminal state $h_{s-1}$.
Reset masks retain their original positions, and $h_{\rm reset}$ stays the global
reset parameter, not the new chunk's incoming state.
\Plate{04-chunk}{Ordinary chunking carries both the value and its differentiation
path. Reset starts a declared new segment; detach keeps values but deliberately
changes the gradient.}{evod26:chunk}
\Listing{chunked}

The reference tests sizes $1,2,3,5$, larger-than-sequence chunks and uneven tails.
It includes an explicit \texttt{detach=True} intervention for a counterexample.
With $a=(1,1)$, $b=(0,0)$, $h_0=1$ and loss $L=h_2$, both ordinary and detached
chunking produce $h_2=1$. Ordinary differentiation gives
$\partial L/\partial b_1=1$; detaching after step one gives zero. Forward parity
alone therefore cannot establish training parity.

Truncated backpropagation may intentionally detach boundaries. It is a different
gradient contract, not a transparent memory optimization. Chapter 27 will separate
gradient accumulation and activation recomputation from such changes.

\section{An independent closed-form oracle}
Unroll Eq.~\ref{evod26:rec} to obtain
\begin{equation}
h_t=\left(\prod_{j=1}^t a_j\right)\odot h_0+
\sum_{k=1}^t\left(\prod_{j=k+1}^t a_j\right)\odot b_k.
\label{evod26:closed}
\end{equation}
An empty product is $\mathbf1$. This formula handles $a_j=0$ without division.
Rewriting every suffix product as a ratio of cumulative products would fail at zero,
precisely where reset semantics become important.

The \texttt{closed\_form} reference evaluates each prefix using products and a
reverse suffix-product loop, rather than the composition operator or forward state
loop. It costs $O(T^2D)$ and exists for verification, not serving.
Tests compare all four schedules across empty, singleton, power-of-two and
non-power-of-two lengths. Hand traces, homogeneous matrices and this product/sum
formula reduce the risk of two implementations sharing the same mistake.

\section{Derive a reverse-mode gradient oracle}
For an assigned linear loss $L=\sum_t w_t^\top h_t$, let
$\lambda_t=\partial L/\partial h_t$. Reverse accumulation gives
\[
\lambda_T=w_T,\qquad
\lambda_t=w_t+a_{t+1}\odot\lambda_{t+1}.
\]
Consequently
\begin{equation}
\frac{\partial L}{\partial b_t}=\lambda_t,\quad
\frac{\partial L}{\partial a_t}=\lambda_t\odot h_{t-1},\quad
\frac{\partial L}{\partial h_0}=a_1\odot\lambda_1.
\label{evod26:adjoint}
\end{equation}
These derivatives are with respect to the provided coefficient tensors. If a
coefficient generator has parameters $\theta$, an end-to-end implementation must
also differentiate through $a(\theta,x),b(\theta,x)$ and compare those gradients.
\Plate{05-gates}{Parity requires independent forward and derivative checks,
including state-boundary interventions. Agreement on one convenient trace is not
a universal stability result.}{evod26:gates}
\Listing{analytic_adjoint}

The reverse oracle implements Eq.~\ref{evod26:adjoint} without asking autograd for
these derivatives. Central finite differences provide another check, including
a zero coefficient. Reset tests compare all participating gradients, including
the reset state, against the literal reset loop. Padding coefficient gradients
are zero because its effective map is independent of those coefficients.

For seed 26, length 17, dimension 5 and float64 coefficients in $[0.8,0.9)$,
the generated result file reports maximum forward discrepancy $2.22045\times10^{-16}$,
gradient discrepancy $4.44089\times10^{-16}$ and manual-adjoint discrepancy zero.
This is a bounded deterministic fixture, not a
general theorem. Long 1024-step bounded-coefficient tests use declared float64
and float32 tolerances; no BF16 or FP16 certification is inferred.

\section{Numerical equivalence has a real counterexample}
Exact-real affine composition is associative. Machine addition and multiplication
are not exactly associative. Consider $a=(1,1,1)$, $b=(10^{16},-10^{16},1)$,
$h_0=0$ in float64. Sequential evaluation gives final state $1$; this scan's
regrouping gives $0$. The absolute difference is one, not a negligible last-bit
artifact at the output scale.

The example exposes cancellation: adding one to a large-magnitude partial sum can
round it away before the opposite large term arrives. It does not refute the
real-arithmetic proof; it refutes an unrestricted numerical-parity claim.
Declare dtype, coefficient/state scales, context lengths and both absolute and
relative tolerances. An allclose pass on contractive random coefficients says
little about a cancellation-heavy or expansive workload.

Long products also express model stability. Values $|a_t|<1$ tend to attenuate
older contributions; values above one can amplify state and gradients.
A gate parameterization may reduce amplification but does not eliminate cancellation,
underflow, accumulation error or data-distribution shift.
The software validates finite inputs; it does not guarantee that every finite
input yields finite intermediates. Production paths need output monitoring and
explicit non-finite failure handling.

\section{Depth, work and measured latency are different quantities}
At offset $2^k$, our scan performs $T-2^k$ pair compositions. For $T>1$,
\[
W_{\rm HS}(T)=\sum_{k=0}^{\lceil\log_2T\rceil-1}(T-2^k),
\qquad
D_{\rm HS}(T)=\lceil\log_2T\rceil.
\]
This is $O(T\log T)$ pair work and $O(\log T)$ idealized stage depth.
It is not a work-efficient Blelloch up-sweep/down-sweep implementation.
For $T=1024$, the reference performs 9217 pair compositions across ten stages;
a literal sequence makes 1024 state updates. A pair composition and a state update
are not identical FLOP units.
\begin{center}\input{\Generated/work.tex}\end{center}

Tensor allocation, concatenation, synchronization and kernel dispatch remain
outside this pair-operation count. With unlimited ideal parallel lanes, lower
dependency depth can help; with finite hardware, extra work and memory traffic
can dominate. No measured speedup is reported. A serious benchmark would specify
hardware, shape, batch, dtype, compilation/fusion, warmup, device synchronization,
peak memory and forward/backward scope.

Mamba's selected implementation discussion emphasizes kernel fusion and avoiding
large intermediates in the memory hierarchy \cite{GuDao2024}. It would be misleading
to assign those properties to this intentionally transparent PyTorch teaching scan.

\section{Serving parity includes ownership}
\Plate{06-life}{State belongs to a logical request, not merely a reusable buffer.
Versioned position, complete state and reset/clone/cancel behavior determine the
meaning of incremental execution.}{evod26:life}
\Term{State ownership} means a request state identifies its model/configuration version, sequence position,
dtype/device and all architecture-specific tensors. For the affine reference that
is $h$; for actual DOGMA it includes native hidden state and memory loci; for
HermonDNA cached attention it includes layer-specific K/V, position and masks.
A common interface may wrap them but cannot replace their semantics.

Clone should preserve independent future evolution through deep copy or a proven
copy-on-write contract. A checkpoint names both position and model version.
Cancellation invalidates the logical owner before buffer reuse. A reset on a new
request is not an optional performance hint. Otherwise serving can leak previous
request information even if the mathematical recurrence is correct.

Chapter 25's Hermon prefix test repeated causal prefixes; it did not implement
KV-cache decoding. This chapter likewise makes no Hermon KV-parity claim.
That requires actual cache construction, correct positions/masks and every-position
logit comparison against the complete causal model.

\section{Failure-oriented verification}
The chapter tests reject swapped composition order, partially updated in-place
scan stages, wrong padding maps, reset-on-padding instructions and detached
boundary gradients. Future-coefficient perturbations must not change earlier
states; independent batch lanes must not interact. Invalid shapes, dtypes,
non-finite inputs and Boolean chunk sizes are rejected.

No one oracle suffices. Matrix products check composition; hand arithmetic checks
a trace; the closed form checks prefix values; the reverse recurrence and finite
differences check derivatives; boundary interventions check lifecycle semantics.
The cancellation fixture remains an expected disagreement, ensuring that tests
do not erase a real numerical limitation merely to make every case pass.

\section{Twenty worked exercises}
\begin{enumerate}
\item Compose $(2,1)$ then $(3,4)$.
\textbf{Answer:} $(6,7)$; reversing gives $(6,9)$.
\item What is the identity?
\textbf{Answer:} $(\mathbf1,\mathbf0)$.
\item Expand a three-step translated term.
\textbf{Answer:} $a_3a_2b_1+a_3b_2+b_3$.
\item Why is a terminal reduction insufficient for training?
\textbf{Answer:} losses may require every position's state.
\item Trace the four-step example from $h_0=2$.
\textbf{Answer:} $(2,4.5,3,-1)$.
\item What does $a_3=0$ erase?
\textbf{Answer:} dependence of $h_3$ on earlier incoming state, while retaining $b_3$.
\item Why not divide cumulative products to compute suffixes?
\textbf{Answer:} zero coefficients make the ratio undefined.
\item Give the pair for reset-before-update.
\textbf{Answer:} $(0,a_t h_{\rm reset}+b_t)$.
\item Give the pair for padding.
\textbf{Answer:} $(1,0)$, not $(0,0)$.
\item Compute states with $a=.5,b=1,h_0=2,h_{\rm reset}=10$ and reset at step two.
\textbf{Answer:} $(2,6,4)$.
\item What if step three is padding?
\textbf{Answer:} its state is $6$.
\item Does resetting cut all gradients to a shared reset parameter?
\textbf{Answer:} no; it cuts earlier-record state dependence, not the new reset contribution.
\item Why preserve the global reset state across chunks?
\textbf{Answer:} using the incoming chunk state as reset changes segment semantics.
\item Give a forward-equal, gradient-unequal detach example.
\textbf{Answer:} two identity-coefficient steps with $L=h_2$ yield
$\partial L/\partial b_1=1$ ordinarily, zero after boundary detach.
\item Derive $\partial L/\partial b_t$.
\textbf{Answer:} the direct state derivative is one, hence $\lambda_t$.
\item Why can causal state-derived coefficients prevent a whole-sequence scan?
\textbf{Answer:} producing coefficients still requires previous recurrent states.
\item Count pair compositions at $T=4$.
\textbf{Answer:} $(4-1)+(4-2)=5$ across two stages.
\item At $T=1024$?
\textbf{Answer:} $10(1024)-(1+\cdots+512)=9217$.
\item Evaluate the cancellation example under the two schedules.
\textbf{Answer:} float64 sequential final $1$, reference scan final $0$.
\item What would justify a production parity claim?
\textbf{Rubric:} actual model/coefficient gradients, complete native state,
resets/padding/positions, representative dtype and length tests, independent
oracles, disclosed tolerances and separately measured systems performance.
\end{enumerate}

\section{Bridge to memory-efficient training}
The next chapter changes how training work and activations are scheduled.
Accumulation must preserve the declared loss weighting; recomputation must preserve
the transition and stochastic behavior; detach changes the gradient unless it is
explicitly part of the objective. The present chapter supplies the value, derivative
and boundary contracts against which those changes must be tested.
